Die abgestrahlte Intensität ist $J = \varepsilon\,\sigma\,T^4$. Ein Stein hat wie die meisten nichtmetallischen Oberflächen eine grosse Emissionszahl (nahe bei 1). Die blanke Aluminiumfolie dagegen hat eine sehr kleine Emissionszahl (poliert $\varepsilon \approx 0.04$, FoTaBe S.~190); sie reflektiert den grössten Teil der Wärmestrahlung des Steins zurück.

Die dünne Folie liegt am Stein an und nimmt durch Wärmeleitung fast seine Temperatur an. Nach aussen strahlt aber nur noch die Folienoberfläche mit ihrer kleinen Emissionszahl. Die Intensität der abgestrahlten Wärme \emph{nimmt deshalb stark ab}, auf einen kleinen Bruchteil (grob ein Zwanzigstel) des ursprünglichen Wertes. Der Stein bleibt länger warm (Prinzip der Rettungsdecke und der Thermoskanne).
